\documentclass[11pt, a4paper]{article}

% --- Pakete ---
\usepackage[english, ngerman]{babel}
\usepackage[utf8]{inputenc}
\usepackage[T1]{fontenc}
\usepackage{lmodern}
\usepackage[left=2.5cm, right=2.5cm, top=2.5cm, bottom=2.5cm]{geometry}
\usepackage[autostyle=true,german=quotes]{csquotes}
\usepackage{graphicx}
\usepackage{xcolor}
\usepackage{enumitem}
\usepackage{tabularx}
\usepackage{booktabs}
\usepackage{tikz}
\usetikzlibrary{positioning,fit,backgrounds,shapes.geometric,arrows.meta}
\usepackage{fancyhdr}
\usepackage{titlesec}
\usepackage{parskip}
\usepackage{microtype}
\usepackage{tocloft}
\usepackage[document]{ragged2e}
\usepackage[backend=biber,style=ieee,urldate=long]{biblatex}
\addbibresource{literatur.bib}
\usepackage{hyperref}
\emergencystretch=0.3\textwidth % lange Wörter ohne Trennung nicht über den Rand laufen lassen
\hyphenpenalty=10000
\exhyphenpenalty=50 % Umbruch nur an vorhandenen Bindestrichen (z.B. Projektmanagement-Werkzeug)

% X-Spalten in tabularx global linksbündig (kein Blocksatz in Tabellen)
\renewcommand{\tabularxcolumn}[1]{>{\raggedright\arraybackslash}p{#1}}

% --- Farben (HKA Corporate) ---
\definecolor{hkared}{RGB}{198, 42, 18}
\definecolor{hkagray}{RGB}{100, 100, 100}

% --- Inhaltsverzeichnis Styling ---
% Überkapitel: rot und fett
\renewcommand{\cftsecfont}{\bfseries\color{hkared}}
\renewcommand{\cftsecpagefont}{\bfseries\color{hkared}}
% Unterkapitel: schwarz (hidelinks im TOC nötig damit hyperref nicht überschreibt)
\renewcommand{\cftsubsecfont}{\normalfont\color{black}}
\renewcommand{\cftsubsecpagefont}{\normalfont\color{black}}
% Spacing: gleiches Spacing auf allen Ebenen
\setlength{\cftbeforesecskip}{8pt}
\setlength{\cftbeforesubsecskip}{8pt}
\setlength{\cftbeforesubsubsecskip}{8pt}
% Abstand nach dem Titel "Inhaltsverzeichnis" vor dem ersten Eintrag
\setlength{\cftaftertoctitleskip}{1.2cm}

% --- Hyperref Setup ---
\hypersetup{
    colorlinks  = true,
    linkcolor   = hkared,
    urlcolor    = hkared,
    citecolor   = hkared,
    pdftitle    = {Expos\'e: Overhead von OpenTelemetry in einer .NET-Microservice-Anwendung},
    pdfauthor   = {Erlind Sejdiu},
    pdfcreator  = {pdflatex},
    pdfproducer = {LaTeX with hyperref}
}

% --- Section Styling ---
\titleformat{\section}{\large\bfseries\color{hkared}}{\thesection}{1em}{}
\titleformat{\subsection}{\normalsize\bfseries\color{hkared}}{\thesubsection}{1em}{}
\titleformat{\subsubsection}{\normalsize\bfseries\color{black}}{\thesubsubsection}{1em}{}
\titlespacing*{\section}{0pt}{1.4em}{0.5em}
\titlespacing*{\subsection}{0pt}{0.8em}{0.3em}
\titlespacing*{\subsubsection}{0pt}{0.6em}{0.2em}

% --- Header/Footer ---
\pagestyle{fancy}
\fancyhf{}
\fancyhead[L]{\small\color{hkagray}Expos\'e -- Projektarbeit}
\fancyhead[R]{\small\color{hkagray}Master Informatik WS~2026/27}
\fancyfoot[C]{\small\color{hkagray}\thepage}
\renewcommand{\headrulewidth}{0.4pt}

% Markierung für Stellen, die noch ergänzt/geprüft werden müssen
\newcommand{\offen}[1]{\textcolor{red}{[#1]}}

% =============================================================================
\begin{document}

% --- Titelseite ---
\begin{titlepage}
    \centering
    \includegraphics[width=0.35\textwidth]{HKA.jpg}

    {\normalsize\color{hkagray} University of Applied Sciences}\\[0.125cm]
    {\small\color{hkagray} Fakult\"at f\"ur Informatik und Wirtschaftsinformatik}

    \vspace{.5cm}
    {\color{hkared}\rule{\textwidth}{1.5pt}}

    \vspace{1.5cm}
    {\fontsize{24}{30}\selectfont\bfseries\color{hkared} Was kostet Observability?}\\[0.6cm]
    {\Large Implementierung einer Microservice-Anwendung mit .NET und Angular
    am Beispiel eines Projektmanagement-Werkzeugs und quantitative Analyse
    des Overheads von OpenTelemetry\par}
    \vspace{.75cm}
    {\LARGE\bfseries Expos\'e}

    \vspace{.5cm}
    {\large Projektarbeit im Master Studiengang Informatik}\\[0.3cm]
    {\large Wintersemester 2026/27}

    \vspace{.5cm}
    {\color{hkared}\rule{0.6\textwidth}{0.5pt}}
    \vspace{1.5cm}

    {\normalsize
    \begin{tabular}{rl}
        \textbf{Autor:}          & Erlind Sejdiu \\
        \textbf{Matrikelnummer:} & 106357 \\[0.5cm]
        \textbf{Korrektor:}      & \offen{Name der betreuenden Person} \\
    \end{tabular}
    }

    \vfill
    {\small\color{hkagray} Stand: \today}
\end{titlepage}

% --- Inhaltsverzeichnis ---
\begingroup
\hypersetup{hidelinks}
\tableofcontents
\endgroup
\thispagestyle{empty}
\newpage
\setcounter{page}{1}

\section{Motivation}

In Microservice-Architekturen durchläuft eine einzelne Nutzeraktion mehrere Dienste, Datenbanken und Nachrichtenkanäle \cite{newman2021microservices}. Tritt ein Fehler auf oder wird eine Anfrage langsam, ist die Ursache oft nicht in einem einzelnen Dienst zu finden, sondern im Zusammenspiel mehrerer Komponenten. \emph{Observability} macht den inneren Zustand eines solchen Systems anhand von Traces, Metriken und Logs nachvollziehbar \cite{majors2022observability}. OpenTelemetry hat sich dafür als herstellerunabhängiger De-facto-Standard etabliert und bietet einheitliche Schnittstellen für viele Sprachen und Plattformen \cite{otel_cncf,otel_docs}.

Observability ist mit Kosten verbunden. Jeder Span, jede Metrik und jeder Log-Eintrag kostet Rechenzeit und Arbeitsspeicher, erzeugt Netzwerkverkehr und belegt Speicherplatz. Zusätzlich kann die Instrumentierung die Antwortzeiten verlängern, die Nutzerinnen und Nutzer direkt wahrnehmen. Wie umfangreich instrumentiert und wie stark gesampelt wird, wird in der Praxis dennoch nach Erfahrung oder den Voreinstellungen der eingesetzten Bibliotheken entschieden. Qualitative Empfehlungen gibt es viele, belastbare quantitative Aussagen darüber, wie groß der Overhead in einer realistischen Anwendung ist und wovon er abhängt, dagegen kaum.

Die Arbeit ist daher zweigeteilt. Im ersten Teil wird ein Projektmanagement-Werkzeug als Microservice-Anwendung mit .NET und Angular entwickelt und mit OpenTelemetry instrumentiert. Im zweiten Teil wird an dieser Anwendung der Einfluss von OpenTelemetry auf Performanz und Speicherbedarf quantitativ gemessen.

\section{Problemstellung}

Nach ISO/IEC 25010 verbessert Observability die \emph{Analysierbarkeit} eines Systems, belastet aber die \emph{Leistungseffizienz} mit ihren Teilmerkmalen \emph{Zeitverhalten} und \emph{Ressourcennutzung} \cite{iso25010}. Um beide Qualitätsmerkmale gegeneinander abwägen zu können, muss der Preis der Observability bekannt sein. Offen ist derzeit:
\begin{enumerate}[nosep]
    \item um wie viel sich Antwortzeiten, Durchsatz einer .NET-Microservice-Anwendung durch OpenTelemetry verschlechtern,
    \item wie viel zusätzlichen Arbeitsspeicher die Dienste benötigen und welches Datenvolumen die Telemetriedaten verursachen,
    \item wie stark Konfigurationsparameter wie die erfassten Signale, die Granularität, das Sampling und die Art des Exports den Overhead beeinflussen.
\end{enumerate}

Eine Antwort auf diese Fragen lässt sich nicht an einem minimalen Beispiel gewinnen. Der Overhead hängt davon ab, wie viele Dienste an einer Anfrage beteiligt sind, wie viel Arbeit jeder Dienst im Verhältnis zur Instrumentierung leistet und wie viele Datenbank- und Netzwerkaufrufe dabei entstehen. Benötigt wird daher eine Anwendung, die typische Anfragen eines realen Systems abbildet und dennoch kontrolliert und wiederholbar gemessen werden kann.

\paragraph{Abgrenzung.} Die Anwendung dient als realistisches Messobjekt und ist kein produktionsreifes Produkt. Nicht betrachtet werden der Vergleich verschiedener Observability-Backends (Granfa oder Datadog). Die Analyse zwischen unterschiedlichen Observability-Backends würde den Rahmen dieser Projektarbeit sprengen.

\section{Stand der Technik}

//

\section{Ziele}

\begin{enumerate}[label=\textbf{\arabic*:}, leftmargin=*, nosep]
    \item Wie stark verändern sich Antwortzeit, maximaler Durchsatz und CPU-Auslastung durch OpenTelemetry?
    \item Wie viel zusätzlichen Arbeitsspeicher benötigen die Dienste, und welches Datenvolumen verursachen die Telemetriedaten?
    \item Welchen Einfluss haben Signale, Granularität, Sampling-Rate und Exportverfahren auf FF1 und FF2?
    \item Wie wirkt sich die Instrumentierung des Angular-Frontends auf Bundle-Größe und Ladezeit aus?
\end{enumerate}

Ergebnis von Teil~1 ist eine lauffähige, in Containern bereitstellbare Anwendung, deren OpenTelemetry-Konfiguration sich ohne Codeänderung umschalten lässt -- bis hin zum Betrieb ganz ohne OpenTelemetry. Ergebnis von Teil~2 sind eine reproduzierbare Messumgebung, die Beantwortung von FF1--FF4 sowie daraus abgeleitete Konfigurationsempfehlungen für Entwicklungsteams.

\section{Vorgehen}\label{sec:vorgehen}

\subsection{Teil 1: Projektmanagement-Werkzeug und Instrumentierung}

Mehrere Personen planen gemeinsam Projekte: Tasks werden angelegt, mit Abhängigkeiten verknüpft und in einer Kanban-, Wochen- und Monatsansicht dargestellt. Beim Verknüpfen wird geprüft, dass keine zyklischen Abhängigkeiten entstehen. Projekte werden über Rollen mit anderen geteilt; Änderungen erscheinen bei allen Beteiligten in Echtzeit.

Abbildung~\ref{fig:architektur} zeigt die geplante Architektur. Das Angular-Frontend kommuniziert ausschließlich mit einem API-Gateway, das Anfragen authentifiziert und an die zuständigen Dienste weiterleitet. Jeder Dienst verwaltet seine Daten in einer eigenen Datenbank. Zur Berechtigungsprüfung rufen Projekt- und Task-Dienst den Nutzer- und Team-Dienst synchron auf. Änderungen werden zusätzlich als Ereignisse über einen Nachrichtenbroker veröffentlicht. Diese Ereignisse verarbeitet der Benachrichtigungsdienst asynchron und sendet daraus Echtzeit-Updates an das Frontend.

\begin{figure}[htbp]
\centering
\resizebox{\textwidth}{!}{% Vereinfachte Architektur des Projektmanagement-Werkzeugs (eingebunden in Expose.tex)
\begin{tikzpicture}[
    font=\small,
    >={Stealth[length=2mm]},
    comp/.style={draw=hkagray, fill=white, rounded corners=3pt, align=center,
                 minimum height=9mm, inner sep=3pt},
    svc/.style={draw=hkared, fill=hkared!12, rounded corners=3pt, align=center,
                minimum width=4.2cm, minimum height=6.5mm},
    db/.style={cylinder, shape border rotate=90, aspect=0.25, draw=hkagray,
               fill=hkagray!15, minimum width=1.3cm, minimum height=8mm, inner sep=1pt},
    box/.style={draw=hkagray!60, fill=hkagray!6, rounded corners=4pt, inner sep=5pt},
    group/.style={draw=hkared, line width=0.8pt, rounded corners=6pt, inner sep=9pt},
    grouptitle/.style={fill=white, inner sep=2pt, font=\small\bfseries\color{hkared}},
    sync/.style={->, draw=black!75},
    event/.style={->, dashed, draw=black!75},
    telemetry/.style={->, dotted, thick, draw=hkagray},
]

% --- Client ---
\node[comp, minimum width=2.4cm] (fe) at (0,2) {Browser\\[1pt]\textbf{Angular-}\\\textbf{Frontend}};

% --- API-Gateway ---
\node[comp, minimum width=2.2cm] (gw) at (3.5,2) {\textbf{API-Gateway}\\[1pt]\footnotesize Authentifizierung};

% --- Dienste (je mit eigener Datenbank) ---
\foreach \name/\label/\y in {user/Nutzer- und Team-Dienst/3.8, proj/Projekt-Dienst/2.6,
                              task/Task-Dienst/1.4, notif/Benachrichtigungsdienst/0.2} {
    \node[svc] (\name) at (7.5,\y) {\label};
    \node[db] (\name db) at (10.6,\y) {\footnotesize DB};
    \draw[sync] (\name.east) -- (\name db.west);
    \begin{scope}[on background layer]
        \node[box, fit=(\name)(\name db)] (\name box) {};
    \end{scope}
}

% --- Nachrichtenbroker ---
\node[comp, minimum width=2.2cm, minimum height=1.4cm] (bus) at (13.8,2) {\textbf{Nachrichten-}\\\textbf{broker}\\[1pt]\footnotesize Ereignisse};

% --- Observability ---
\node[comp, minimum width=2.6cm] (col) at (6.2,-2.5) {\textbf{Collector}};
\node[comp, minimum width=3.4cm] (obs) at (10.6,-2.5) {\textbf{Analyse-Backend}\\[1pt]\footnotesize Traces, Metriken, Logs};

% --- Gruppen ---
\begin{scope}[on background layer]
    \node[group, fit=(fe)] (gclient) {};
    \node[group, fit=(gw)(userbox)(notifbox)(bus)] (gbackend) {};
    \node[group, draw=hkagray, fit=(col)(obs)] (gobs) {};
\end{scope}
\node[grouptitle] at (gclient.north) {Client};
\node[grouptitle] at (gbackend.north) {Backend (Container)};
\node[grouptitle, text=hkagray] at (gobs.north) {Observability};

% --- Synchrone Aufrufe ---
\draw[sync] (fe.east) -- (gw.west);
\draw (gw.east) -- (5.0,2);
\draw (5.0,3.8) -- (5.0,0.2);
\foreach \name/\y in {userbox/3.8, projbox/2.6, taskbox/1.4, notifbox/0.2}
    \draw[sync] (5.0,\y) -- (\name.west |- 0,\y);

% --- Ereignisse ---
\draw[dashed, draw=black!75] (projbox.east |- 0,2.6) -| (12.0,0.2);
\draw[dashed, draw=black!75] (taskbox.east |- 0,1.4) -- (12.0,1.4);
\draw[event] (12.0,2) -- (bus.west);
\draw[event] (12.0,0.2) -- (notifbox.east |- 0,0.2);
\draw[event] ([xshift=-4mm]notifbox.south) -- ++(0,-0.95) -| node[pos=0.25, below, inner sep=2pt, font=\footnotesize] {Echtzeit-Updates} ([xshift=-5mm]gclient.south);

% --- Telemetrie ---
\draw[telemetry] ([xshift=5mm]gclient.south) |- (col.west);
\draw[telemetry] (col.north |- gbackend.south) -- (col.north);
\draw[telemetry] (col.east) -- node[above, font=\footnotesize, text=hkagray] {OTLP} (obs.west);

\end{tikzpicture}
}
\caption{Geplante Architektur des Projektmanagement-Werkzeugs (vereinfacht). Durchgezogene Pfeile: synchrone Aufrufe; gestrichelt: Ereignisse; gepunktet: Telemetriedaten.}
\label{fig:architektur}
\end{figure}

Die Domäne enthält damit unterschiedliche Anfragepfade, die für die Messung relevant sind: große Lesezugriffe (Kalenderansichten, Board), kurze Schreibzugriffe (Task verschieben), rechenintensive Logik (Zyklusprüfung), synchrone Aufrufketten über mehrere Dienste und asynchrone Verarbeitung.

Instrumentiert wird zunächst automatisch, sodass eingehende und ausgehende Anfragen, Datenbankzugriffe und Nachrichten ohne eigenen Code erfasst werden. Ergänzt wird dies um manuelle Spans und Metriken für fachliche Vorgänge wie die Zyklusprüfung. Das Frontend gibt den Trace-Kontext an das Backend weiter, sodass eine Nutzeraktion vom Klick bis zur Datenbank als ein Trace sichtbar wird. Alle Dienste senden ihre Daten an einen gemeinsamen Collector, der sie an ein Analyse-Backend weiterleitet. Welche Signale erfasst werden, wie fein instrumentiert und wie stark gesampelt wird, lässt sich über Konfiguration umschalten. Wichtige Entwurfsentscheidungen werden als Architecture Decision Records dokumentiert \cite{nygard2011adr}.

\subsection{Teil 2: Messung}

\paragraph{Versuchsdesign.} Gemessen wird gegen eine Basislinie ohne OpenTelemetry. Ausgehend von einer Standardkonfiguration wird jeweils ein Faktor variiert, um seinen Einfluss isoliert betrachten zu können:

\begin{center}
\small
\begin{tabularx}{\textwidth}{@{}l X@{}}
\toprule
\textbf{Faktor} & \textbf{Stufen} \\
\midrule
Instrumentierung & keine; SDK ohne Export; automatisch; zusätzlich manuell (grob/fein) \\
Signale & Traces; zusätzlich Metriken; zusätzlich Logs \\
Sampling & 0\,\%, 1\,\%, 10\,\%, 100\,\% \\
Export & verschiedene Übertragungsprotokolle; gebündelter vs.\ sofortiger Export \\
Last & 25\,\%, 50\,\%, 75\,\% des maximalen Durchsatzes der Basislinie \\
\bottomrule
\end{tabularx}
\end{center}

\paragraph{Messgrößen.} Erfasst werden Antwortzeiten und Durchsatz aus Sicht eines Lastgenerators, CPU und Arbeitsspeicher je Dienst einschließlich Speicherallokationen und Garbage Collection, das Datenvolumen zum Collector und der belegte Speicherplatz im Backend. Im Frontend werden Bundle-Größe und Ladezeit anhand der Core Web Vitals bewertet \cite{webvitals}.

\paragraph{Messumgebung.} Der Lastgenerator führt gewichtete Szenarien aus, die typische Nutzung abbilden: Board laden, Wochen- und Monatsansicht öffnen, Tasks anlegen, verschieben und verknüpfen. Grundlage ist ein fester, generierter Datenbestand, damit alle Konfigurationen unter gleichen Bedingungen gemessen werden. Ressourcen werden mit Werkzeugen erfasst, die unabhängig von OpenTelemetry arbeiten, damit die Messung das Ergebnis nicht verfälscht. Einzelne Operationen werden zusätzlich isoliert per Micro-Benchmark gemessen. Alle Container erhalten feste Ressourcenlimits, jede Konfiguration wird nach einer Aufwärmphase mehrfach und automatisiert gemessen. \offen{Messrechner klären.}

\paragraph{Auswertung.} Angegeben werden Perzentile mit Konfidenzintervallen und der relative Overhead zur Basislinie; Unterschiede werden mit dem Mann-Whitney-U-Test auf Signifikanz geprüft \cite{georges2007rigorous,kalibera2013rigorous}. Die Ergebnisse werden mit bestehenden Studien verglichen \cite{reichelt2023overhead}. Abschließend werden die Übertragbarkeit auf reale Systeme sowie der Einfluss von Messumgebung und Softwareversionen diskutiert.

\section{Zeitplan}

\begin{center}
\small
\begin{tabularx}{\textwidth}{@{}l X X@{}}
\toprule
\textbf{Phase} & \textbf{Ziel} & \textbf{Artefakt} \\
\midrule
1 Einarbeitung & OpenTelemetry, Messmethodik & Literaturübersicht \\
2 Implementierung & Dienste und Angular-Frontend & Lauffähige Anwendung, ADRs \\
3 Instrumentierung & Backend, Frontend, Collector & Umschaltbare Konfiguration \\
4 Messumgebung & Lastszenarien, Automatisierung, Pilotmessungen & Mess- und Auswertungsskripte \\
5 Messungen & Alle Konfigurationen messen & Rohdaten \\
6 Auswertung & Statistik, Empfehlungen & Diagramme, Tabellen \\
7 Ausarbeitung & Schriftliche Dokumentation & Ausarbeitung, Code und Daten im Repository \\
\bottomrule
\end{tabularx}
\end{center}

\printbibliography[heading=bibintoc, title={Literaturverzeichnis}]

\end{document}
